\documentclass[12pt,a4paper]{article}

% ---- Language & Fonts ----
\usepackage[T2A]{fontenc}
\usepackage[utf8]{inputenc}
\usepackage[russian,english]{babel}
\selectlanguage{russian}
\usepackage{amsmath,amssymb,amsfonts,mathtools}
\usepackage{bm}
\usepackage{graphicx}
\usepackage{booktabs}
\usepackage{siunitx}
\usepackage{hyperref}
\usepackage[margin=2.5cm]{geometry}
\usepackage{natbib}
\bibliographystyle{unsrtnat}
\usepackage{algorithm}
\usepackage{algorithmic}
\usepackage{xcolor}
\usepackage{enumitem}
\usepackage{caption}
\usepackage{subcaption}

% ---- Custom commands ----
\newcommand{\Eone}{E_{1}}
\newcommand{\vu}{\mathbf{u}}
\newcommand{\vz}{\mathbf{z}}
\newcommand{\va}{\mathbf{a}}
\newcommand{\vb}{\mathbf{b}}
\newcommand{\vr}{\mathbf{r}}
\newcommand{\mK}{\mathbf{K}}
\newcommand{\mA}{\mathbf{A}}
\newcommand{\mL}{\mathbf{L}}
\newcommand{\mD}{\mathbf{D}}
\newcommand{\mR}{\mathbf{R}}
\newcommand{\mI}{\mathbf{I}}
\newcommand{\mW}{\mathbf{W}}
\newcommand{\mC}{\mathbf{C}}

\title{
  Инверсия интегрального уравнения Фредгольма первого рода\newline
  для восстановления вертикального профиля объёмной активности\newline
  радионуклидов в почве с учётом миграционной химии\newline
  и методов геофизической регуляризации
}

\author{
  Автор~[1] \\
  \smallskip
   \small [1] Институт проблем безопасного развития атомной энергетики РАН
}

\date{2026}

\begin{document}

\maketitle

% =====================================================================
\section*{Аннотация}
\addcontentsline{toc}{section}{Аннотация}

Представлена программная реализация и теоретическое обоснование комплекса методов для
восстановления вертикального профиля объёмной активности радионуклидов $a(z)$
по результатам 	extit{in-situ} гамма-спектрометрии.
Задача сводится к решению интегрального уравнения Фредгольма первого рода с ядром,
основанным на формуле Бека--де Планка (1968) и выражающемся через
интегральную экспоненциальную функцию $\Eone$.

Научная новизна работы состоит в систематическом переносе современных методов
геофизической инверсии (1996--2025) в область гамма-спектрометрического
зондирования почвы.
Впервые для данной задачи реализованы и объединены в единый программный конвейер:
итеративные методы (CGLS, метод наименьших квадратов с сопряжёнными градиентами,
с ранней остановкой по L-кривой; рандомизированный метод Качмаржа),
компактная инверсия по Порнягину--Жданову (IRLS, минимальная градиентная поддержка),
регуляризация полной вариации (TV/ADMM),
разреженная L1-инверсия (FISTA с ускорением Нестерова),
ансамблевая инверсия Калмана (EKI) с квантованием неопределённости,
а также ансамблевый мультиметодовый подход с отбором моделей по критериям Акаике (AIC)
и Байесовскому информационному критерию (BIC).

Дополнительно развит блок миграционной химии, включающий уравнение адвекции--дисперсии
с кинетической сорбцией по двухузловой модели ван Генухтена--Вагенета,
pH-зависимый коэффициент распределения, изотерму Фрейндлиха,
конкурентный ионный обмен (Cs$^+$--K$^+$--NH$_4^+$, Sr$^{2+}$--Ca$^{2+}$),
Eh-зависимую сорбцию урана, цепочки распада Бейтмена и численное решение
многослойного ADE методом Кранка--Николсона.

Точность и устойчивость методов демонстрируются на численных экспериментах
с синтетическими данными для Eu-152 (три гамма-линии, две высоты детектора).

\medskip
\noindent\textbf{Ключевые слова:}
интегральное уравнение Фредгольма, гамма-спектрометрия, инверсия глубинного
профиля, регуляризация Тихонова, полная вариация, компактная инверсия,
ансамблевая инверсия Калмана, уравнение адвекции--дисперсии, коэффициент
распределения, кинетическая сорбция.

% =====================================================================
\tableofcontents
\newpage

% =====================================================================
\section{Введение}
\label{sec:intro}

Определение вертикального распределения радионуклидов в почве является
фундаментальной задачей радиационной разведки, мониторинга загрязнений
и оценки дозовых нагрузок на население.
Полевая 	extit{in-situ} гамма-спектрометрия позволяет получить информацию
о поверхностной активности и, при определённых условиях, о глубинном
распределении радионуклидов без отбора проб~\citep{Beck1968, Zombori1992, Tyler2008}.

Метод основан на том, что угловое распределение гамма-излучения от
точечного источника на глубине $z$ несёт информацию об этой глубине:
поверхностные источники дают более изотропное поле, тогда как глубокие
источники---более коллимированное. Детектор, поднятый на высоту $h$ над
поверхностью, регистрирует поток, пропорциональный интегралу
от глубины $0$ до $z_{\max}$, взвешенному ядром $k(z)$,
что формирует интегральное уравнение Фредгольма первого рода~\citep{Beck1968}.

Задача инверсии такого уравнения является некорректно поставленной
в смысле Адамара: малые возмущения данных (статистический шум счёта,
погрешность калибровки эффективности) приводят к неограниченным
колебаниям решения~\citep{Engl1996, Hansen1998}.
Классический подход к решению основан на регуляризации Тихонова
с выбором параметра регуляризации по обобщённому критерию
перекрёстной проверки (GCV)~\citep{Golub1979}, L-кривой~\citep{Hansen1992}
или принципу невязки Морозова~\citep{Morozov1966}.

В последние десятилетия в геофизике (гравиразведка, магниторазведка,
электроразведка, сейсмика) был разработан обширный арсенал методов
регуляризованной инверсии, адаптированных к специфической структуре
ядрForward-оператора~\citep{LiOldenburg1996, LiOldenburg1998, Zhdanov2002,
Portniaguine1999, Vatankhah2018, Chen2024}.
Эти методы учитывают глубинное затухание чувствительности,
способствуют получению компактных решений с резкими границами,
количественно оценивают неопределённость.

\textbf{Цель настоящей работы}---перенести эти методы в область
гамма-спектрометрического зондирования почвы, объединить их с
моделями миграционной химии радионуклидов и реализовать в виде
открытого программного обеспечения.

Основные элементы научной новизны, фиксируемые в данной работе:

\begin{enumerate}[leftmargin=*,itemsep=3pt]
  \item Применение методов компактной инверсии (MGS/MS по Порнягину--Жданову)
        к задаче восстановления профиля $a(z)$, обеспечивающих
        блочно-слоистые решения с резкими границами раздела.
  \item Реализация регуляризации полной вариации (TV) через схему
        ADMM для сохранения резких границ загрязнённых слоёв.
  \item Адаптация метода FISTA (ускоренный проксимальный градиент) с L1-штрафом
        для получения разреженных компактных профилей активности.
  \item Внедрение CGLS с автоматической остановкой по L-кривой
        (аналог выбору $\alpha$ в Тихонове) и рандомизированного метода Качмаржа.
  \item Ансамблевая инверсия Калмана (EKI) с количественной оценкой
        апостериорной неопределённости без использования MCMC.
  \item Мультиметодовый ансамбль с отбором лучшей модели по AIC/BIC.
  \item Интеграция миграционной химии: pH-зависимый $K_d$,
        изотерма Фрейндлиха, конкурентный ионный обмен, двухузловая
        кинетическая сорбция, Eh-зависимая сорбция урана.
  \item Совместная инверсия нескольких нуклидов с общей глубинной сеткой.
  \item Расширенный набор критериев выбора параметра регуляризации:
        квази-оптимальность, NCP (нормализованный кумулятивный периодограммный),
        SNR-критерий, взвешенный GCV для пуассоновского шума.
  \item Геофизические диагностические инструменты: матрица разрешающей
        способности, глубина исследования (DOI), функция расплывания,
        шахматный тест, информационное содержание.
\end{enumerate}


% =====================================================================
\section{Математическая постановка задачи}
\label{sec:problem}

\subsection{Уравнение Фредгольма для гамма-зондирования}
\label{sub:fredholm}

Детектор гамма-излучения, расположенный на высоте $h$ над поверхностью
почвы, регистрирует счётную скорость в $i$-й полной поглощённой пике,
связанную с вертикальным профилем объёмной активности $a(z)$
[Бк/м$^3$] интегральным уравнением Фредгольма первого рода:

\begin{equation}
  C_i = \int_{0}^{z_{\max}} k_i(z)\, a(z)\, dz,\quad i = 1, \ldots, m,
  \label{eq:fredholm}
\end{equation}

где $m$---число измерительных каналов (пик $\times$ высота),
а ядро $k_i(z)$ описывает вероятность регистрации фотона
от точечного изотропного источника на глубине $z$.

\subsection{Ядро: формула Бека--де Планка}
\label{sub:kernel}

При допущении об изотропном детекторе и отсутствии рассеяния
(нет фактора накопления) ядро выражается через интегральную
экспоненциальную функцию $\Eone(x) = \int_x^\infty e^{-t}/t\, dt$~\citep{Beck1968}:

\begin{equation}
  k_i(z) = \frac{y_i\, \varepsilon}{2}\, \Eone\!bigl(\mu_{\mathrm{air}}\, h + \mu_{\mathrm{soil}}\, z\bigr),
  \label{eq:kernel_analytic}
\end{equation}

где $y_i$---выход гамма-квантов на распад в $i$-й линии,
$\varepsilon$---полнотелесная эффективность детектора,
$\mu_{\mathrm{soil}}$ и $\mu_{\mathrm{air}}$---линейные коэффициенты
ослабления в почве и воздухе соответственно.

В общем случае (с учётом угловой зависимости эффективности $g(\theta)$
и фактора накопления $B(\tau)$) ядро вычисляется численным интегрированием:

\begin{equation}
  k_i(z) = \frac{y_i\, \varepsilon}{2}
  \int_{0}^{\pi/2} \tan\theta\; g_i(\theta)\; B(\tau)\; e^{-\tau}\, d\theta,\quad
  \tau = \frac{\mu_{\mathrm{air}}\, h + \mu_{\mathrm{soil}}\, z}{\cos\theta}.
  \label{eq:kernel_full}
\end{equation}

Интегрирование проводится по переменной $u = 1/\cos\theta$ на
логарифмической сетке методом Гаусса--Лежандра.

\subsection{Расширенные модели ядра}
\label{sub:extended_kernels}

Для реальных условий измерений реализованы три расширенные модели ядра:

\begin{itemize}[leftmargin=*,itemsep=2pt]
  \item \textbf{Латерально-распространённый источник} ($\texttt{kernel\_lateral}$):
        учитывает вклад гамма-фотонов от точек, смещённых
        горизонтально от оси детектора в пределах кругового
        поля зрения радиусом $r_{\mathrm{fov}}$.
        Для точки $(r, z)$ оптическая толщина:
        $\tau = \mu_{\mathrm{soil}}\sqrt{z^2 + r^2} + \mu_{\mathrm{air}}\sqrt{h^2 + r^2}$.
        Необходимо для воздушных и широкополосных наземных съёмок~\citep{Tyler2008}.

  \item \textbf{Коллимированный детектор} ($\texttt{kernel\_collimated}$):
        ядро ограничено конусом коллиматора с половинным углом
        $\theta_c$;
        интегрирование проводится до $u_{\max} = 1/\cos\theta_c$
        с плавным окном на основе $\tanh$.

  \item \textbf{Многослойная почва} ($\texttt{kernel\_multilayer}$):
        каждая слой имеет свой $\mu_{\mathrm{soil}}$;
        оптическая толщина на глубине $z_j$ вычисляется
        с учётом свойств всех вышележащих слоёв.
        Важно для реальных почв, где органический слой
        (низкая плотность) залегает над минеральным.
\end{itemize}

\subsection{Дискретизация}
\label{sub:discretization}

Разбив интервал $[0, z_{\max}]$ на $n$ ячеек шириной $\Delta z_j$,
получаем линейную систему:

\begin{equation}
  \mK\, \va = \vb,\quad
  K_{ij} = k_i(z_j)\, \Delta z_j,\quad
  b_i = C_i,\quad
  \mK \in \mathbb{R}^{m \times n}.
  \label{eq:linear}
\end{equation}

Система является некорректно обусловленной: отношение максимального
и минимального сингулярных чисел $\kappa(\mK)$ может достигать $10^6$--$10^8$.
Прямое обращение (МНК) приводит к неустойчивому решению с
 нефизичными осцилляциями.


% =====================================================================
\section{Методы регуляризованной инверсии}
\label{sec:solvers}

Все решатели реализуют принцип:\ «\textit{решить минимальную невязку}
\textit{при приемлемой гладкости (или компактности) решения}\».
Формально:

\begin{equation}
  \hat{\va} = \arg\min_{\va \geq 0}
  \left\{
  \left\|\frac{\mK\va - \vb}{\bm{\sigma}}\right\|^2
  + \alpha\, \Omega(\va)\right\},
  \label{eq:general}
\end{equation}

где $\bm{\sigma}$---стандартные отклонения данных (пуассоновские:
$\sigma_i = \sqrt{C_i}$), $\alpha$---параметр регуляризации,
$\Omega(\va)$---стабилизирующий функционал.

\subsection{Классические методы}
\label{sub:classical}

\subsubsection{Регуляризация Тихонова}

При $\Omega(\va) = \|\mL(\va - \va_0)\|^2$ задача сводится к NNLS
(метод наименьших квадратов с неотрицательными ограничениями)
через расширённую систему~\citep{Tikhonov1977, Hansen1998}:

\begin{equation}
  \begin{pmatrix} \mA \\ \sqrt{\alpha}\,\mL \end{pmatrix} \va =
  \begin{pmatrix} \vb \\ \sqrt{\alpha}\,\mL\,\va_0 \end{pmatrix},
  \quad \va \geq 0,
\end{equation}

где $\mA = \mK / \sigma$---взвешенная матрица системы,
$\mL$---оператор регуляризации (единичная или конечно-разностная
матрица первого или второго порядка).
Реализовано через \texttt{scipy.optimize.lsq\_linear}.

\subsubsection{ТСVD---усечённое сингулярное разложение}

Сингулярное разложение $\mA = \mU\,\mathrm{diag}(s)\,\mV^T$
в пространстве нормализованных по глубине столбцов;
отбрасываются сингулярные числа ниже порога $s[0] \times \texttt{rel\_cutoff}$.
Нормализация столбцов по методу Ли--Олденбурга~\citep{LiOldenburg1996}
применяется для выравнивания чувствительности по глубине.

\subsubsection{Итерационные методы: Ландвебер и Чиммино}

Метод Ландвебера (метод наискорейшего спуска для нормальных уравнений):
$\va^{k+1} = \va^k + \omega\, \mA^T(\vb - \mA\va^k)$, $\va \geq 0$,
с шагом $\omega < 2/s_{\max}^2$ и остановкой по принципу невязки.
Проявляет свойство полусходимости: ошибка сначала убывает,
затем растёт из-за усиления шума.

Метод Чиммино (одновременный Качмарц)---строчный метод,
обновляющий решение по всем уравнениям одновременно:
$\va^{k+1} = \va^k + \frac{1}{m} \sum_{i=1}^m
  \frac{r_i}{\|\mA_{i\cdot}\|^2}\,\mA_{i\cdot}^T$.

\subsection{Методы геофизической инверсии}
\label{sub:geophysics_solvers}

\subsubsection{CGLS с L-кривой}
\label{sub:cgls}

Метод сопряжённых градиентов для нормальных уравнений
(Conjugate Gradient Least Squares) решает $\mA^T\mA\,\va = \mA^T\vb$
итеративно, при этом число итераций $k$ играет роль параметра
регуляризации~\citep{Chen2024}.
Для автоматической остановки реализован метод L-кривой
(\texttt{cgls\_lcurve}): на каждой итерации записывается невязка
$\|\vr\|^2$ и норма решения $\|\va\|$; оптимальная итерация
определяется максимумом кривизны Менгера (аналог угла L-кривой
для Тихонова).

\textbf{Новизна:} комбинация CGLS с L-кривой для задачи
гамма-спектрометрического зондирования ранее не применялась.
Это позволяет получить решение, сравнимое по качеству с Тихоновым,
но без необходимости явного вычисления SVD.

\subsubsection{Метод Качмаржа (ART) с рандомизацией}

Последовательная проекция на гиперплоскости уравнений:
$\va^{k+1} = \va^k + (b_i - \mA_{i\cdot}\va^k)/\|\mA_{i\cdot}\|^2 \cdot \mA_{i\cdot}$.
Рандомизация порядка строк (Strohmer, Vershynin, 2009) обеспечивает
экспоненциальную сходимость для систем с нормированными строками.

\subsubsection{FISTA: L1-разреженная инверсия}
\label{sub:fista}

Метод быстрого итеративного сжимающего порога (Fast Iterative
Shrinkage-Thresholding Algorithm) решает~\citep{BeckTeboulle2009}:

\begin{equation}
  \hat{\va} = \arg\min_{\va \geq 0}
  \left\{\frac{1}{2}\|\mA\va - \vb\|^2 + \alpha\,\|\va\|_1\right\}.
\end{equation}

L1-штраф промотирует разреженность: решение концентрируется
в нескольких тонких слоях, что физически соответствует
локализованным загрязнённым зонам.
Ускорение Нестерова обеспечивает скорость сходимости $O(1/k^2)$
вместо $O(1/k)$ для обычного ISTA.
Константа Липшица $L = s_{\max}^2$ вычисляется по максимальному
сингулярному числу $\mA$.

\textbf{Новизна:} применение FISTA для инверсии профиля $a(z)$
в гамма-спектрометрии является оригинальным.
Метод особенно эффективен для профилей с резкими пиками
активности (локальные загрязнения).

\subsubsection{TV/ADMM: регуляризация полной вариации}
\label{sub:tv}

Полная вариация $\mathrm{TV}(\va) = \sum_j |(\mD\va)_j|$
штрафует разности соседних ячеек, сохраняя при этом
резкие границы (в отличие от L2-регуляризации Тихонова,
которая сглаживает границы)~\citep{Vatankhah2018}.

Задача решается методом штрафных функций ADMM
(Alternating Direction Method of Multipliers):

\begin{equation}
  \begin{aligned}
    &\min_{\va,\vz}\; \frac{1}{2}\|\mA\va - \vb\|^2 + \alpha\,\|\vz\|_1
    \quad\text{при}\quad \mD\va = \vz,\n    &\text{m-шаг: } (\mA^T\mA + \rho\,\mD^T\mD)\,\va = \mA^T\vb + \rho\,\mD^T(\vz - \vu),\n    &\text{z-шаг: } \vz = S_{\alpha/\rho}(\mD\va + \vu),\n    &\text{u-шаг: } \vu = \vu + \mD\va - \vz,
  \end{aligned}
  \label{eq:admm}
\end{equation}

где $S_{\tau}$---оператор мягкого порога,
$\rho$---штрафной параметр,
$\vu$---двойственная переменная.
Система на m-шаге решается через заранее вычисленное
разложение Холецкого матрицы $\mA^T\mA + \rho\,\mD^T\mD$.
Реализованы как анизотропный ($\sum|D_i|$), так и изотропный
($\sum\sqrt{D_i^2 + \varepsilon^2}$) варианты.

\textbf{Новизна:} TV-регуляризация через ADMM для гамма-спектрометрической
инверсии профиля $a(z)$ реализована впервые.
Метод принципиально важен для profiles с резкими
границами слоёв (например, граница <<органический--минеральный>>
или <<загрязнённый--чистый>> слой).

\subsubsection{Компактная инверсия: IRLS по Порнягину--Жданову}
\label{sub:focusing}

Метод минимальной градиентной поддержки (MGS, Minimum Gradient
Support)~\citep{Portniaguine1999, Zhdanov2002} минимизирует функционал:

\begin{equation}
  \Omega_{\mathrm{MGS}}(\va) = \sum_j \frac{(\mD\va)_j^2}{(\mD\va)_j^2 + e^2},
  \label{eq:mgs}
\end{equation}

где $e$---параметр фокусировки.
Функционал приближённо равен числу ненулевых градиентов,
поэтому его минимизация приводит к блочным решениям с плоскими
внутренними областями и резкими границами.

Альтернатива---минимальная поддержка (MS, Minimum Support):

\begin{equation}
  \Omega_{\mathrm{MS}}(\va) = \sum_j \frac{a_j^2}{a_j^2 + e^2},
  \label{eq:ms}
\end{equation}

которая штрафует число ненулевых элементов, давая разреженные решения.

Оба функционала решаются методом итеративно перевзвешенных
наименьших квадратов (IRLS): на каждой итерации $k$
функционал линеаризуется, и решается взвешенная задача Тихонова
с диагональной матрицей весов $\mW_k = \mathrm{diag}(\mD\va_k^2 + e^2)$.
Параметр $e$ прогрессивно уменьшается ($e_{k+1} = 0.7\,e_k$)
для последовательного уточнения границ~\citep{Zhdanov2002}.

\textbf{Новизна:} компактная инверсия MGS/MS по Порнягину--Жданову,
широко применяемая в гравиразведке и магниторазведке,
впервые адаптирована для инверсии профиля $a(z)$
в гамма-спектрометрии почвы.

\subsection{Выбор параметра регуляризации}
\label{sub:criteria}

Реализован расширенный набор критериев:

\begin{itemize}[leftmargin=*,itemsep=2pt]
  \item \textbf{GCV} (Generalized Cross-Validation)~\citep{Golub1979}:
        $\mathrm{GCV}(\alpha) = n\,\|\mA\va_\alpha - \vb\|^2 / [\mathrm{tr}(\mI - \mH_\alpha)]^2$,
        след матрицы hat-матрицы оценивается стохастическим
        оценщиком Хатчинсона ($n_{\mathrm{probe}} = 24$ вектора).

  \item \textbf{L-кривая}~\citep{Hansen1992}:
        максимум кривизны Менгера в лог-лог пространстве
        $(\|\mL\va\|,\|\vr\|)$.

  \item \textbf{Принцип невязки Морозова}~\citep{Morozov1966}:
        $\|\vr\|^2 = m$ (для единичных весов);
        поиск методом Брента на лог-сетки.

  \item \textbf{Квази-оптимальность}~\citep{Hochstenbach2015}:
        минимизация шумовой компоненты SVD; не требует
        оценки уровня шума.

  \item \textbf{NCP} (Normalized Cumulative Periodogram):
        критерий белизны остатков (KS-тест);
        выбирается $\alpha$, при котором остатки ближе всего
        к белому шуму.

  \item \textbf{SNR-критерий}: максимизация отношения
        <<сигнал/шум>> в решении.

  \item \textbf{Взвешенный GCV}: модификация GCV для
        гетероскедастического (пуассоновского) шума.

  \item \textbf{K-fold перекрёстная проверка} ($\texttt{crossval\_alpha}$):
        шумонезависимая альтернатива GCV.

  \item \textbf{L-кривая для итеративных методов}
        ($\texttt{lcurve\_corner\_iter}$): адаптация L-кривой
        для CGLS и Ландвебера, где параметром является
        номер итерации.
\end{itemize}


% =====================================================================
\section{Байесовские методы}
\label{sec:bayesian}

\subsection{Ансамблевая инверсия Калмана (EKI)}
\label{sub:eki}

Метод ансамблевой инверсии Калмана~\citep{Iglesias2013, Laby2025, Tso2024}
аппроксимирует байесовское апостериорное распределение
$p(\va\,|\,\vb)$ без использования MCMC.
Ансамбль из $N_e$ реализаций модели обновляется
по уравнениям фильтра Калмана:

\begin{enumerate}[leftmargin=*,itemsep=2pt]
  \item \textbf{Forward:} $\mathbf{d}_j^k = \mK\,\va_j^k$ для каждого $j = 1,\ldots,N_e$.
  \item \textbf{Статистика ансамбля:}
        $\mC_{md} = \frac{1}{N_e - 1}\mA_{\mathrm{ens}}^T\,\mD_{\mathrm{ens}}$,
        $\mC_d = \frac{1}{N_e - 1}\mD_{\mathrm{ens}}^T\,\mD_{\mathrm{ens}} + \mC_{dd}$.
  \item \textbf{Обновление Калмана:}
        $\va_j^{k+1} = \va_j^k + \mC_{md}\,\mC_d^{-1}(\vb + \bm{\varepsilon}_j - \mathbf{d}_j^k)$,
        где $\bm{\varepsilon}_j \sim \mathcal{N}(0, \mC_{dd})$---инфляционная
        добавка, предотвращающая коллапс ансамбля.
\end{enumerate}

Регуляризованный вариант (Laby et al., 2025) добавляет $\alpha\mI$
к $\mC_d$ и $\mC_{md}$ для стабильности.

Апостериорное среднее $\bar{\va}$ и стандартное отклонение $\sigma_{\mathrm{post}}$
вычисляются directly из финального ансамбля.

\textbf{Новизна:} EKI для инверсии профиля радионуклидов
по гамма-спектрам реализована впервые.
Это обеспечивает количественную оценку неопределённости
за секунды (в отличие от MCMC, требующего часов)
при сопоставимой точности.

\subsection{Лапласовское приближение MAP}
\label{sub:laplace}

Для гауссовского априорного $p(\va) \sim \mathcal{N}(\va_0, (\alpha\mL^T\mL)^{-1})$
и гауссовского правдоподобия $p(\vb|\va) \sim \mathcal{N}(\mK\va, \mathrm{diag}(\sigma^2))$
MAP-оценка эквивалентна решению Тихонова:

\begin{equation}
  \va_{\mathrm{MAP}} = (\mA^T\mA + \alpha\mL^T\mL)^{-1}\,\mA^T\vb,
\end{equation}

а апостериорная ковариация:

\begin{equation}
  \mathrm{Cov}_{\mathrm{post}} = (\mA^T\mA + \alpha\mL^T\mL)^{-1}.
  \label{eq:post_cov}
\end{equation}

Диагональ $\mathrm{Cov}_{\mathrm{post}}$ даёт $1\sigma$-неопределённость
в каждой ячейке глубины.

Реализован также RBF-гауссовский процесс как априорная ковариация:
$C_{ij} = \sigma_f^2 \exp(-|z_i - z_j|^2 / 2l^2) + \sigma_n^2 \delta_{ij}$,
что кодирует априорное представление о гладкости профиля
с корреляционной длиной $l$~\citep{Hasan2023}.


% =====================================================================
\section{Глубинное взвешивание и совместная инверсия}
\label{sec:geophysics}

\subsection{Функции глубинного взвешивания}
\label{sub:depth_weight}

Естественное экспоненциальное затухание ядра с глубиной приводит
к тому, что без дополнительного взвешивания регуляризация
сильно сглаживает решения на глубине.
Для компенсации реализованы четыре метода:

\begin{itemize}[leftmargin=*,itemsep=2pt]
  \item \textbf{Ли--Олденбург}~\citep{LiOldenburg1996}:
        $S_i = 1/\|\mK_{\cdot i}\|$; нормализация столбцов матрицы $\mK$.
  \item \textbf{Степенное} ($w(z) = (z/z_{\mathrm{ref}} + 1)^{-\beta/2}$):
        адаптация метода из 3D-гравиразведки~\citep{LiOldenburg1998};
        для гамма-ядер типичны $\beta \in [1, 3]$.
  \item \textbf{Логарифмическое} ($w(z) = 1/\ln(1 + z/z_0)$):
        более мягкое взвешивание.
  \item \textbf{Адаптивное}: комбинация нормы столбцов, степенного
        и гладкостного корректиров; корректирует глубины
        с плохой чувствительностью.
\end{itemize}

Глубинное взвешивание включается в регуляризацию как
$\Omega(\va) = \|\mW^{1/2}\mL(\va - \va_0)\|^2$.

\subsection{Совместная инверсия нескольких нуклидов}
\label{sub:joint}

Для одновременного определения профилей $a_1(z), a_2(z), \ldots$
нескольких радионуклидов строится блочно-диагональная система:

\begin{equation}
  \begin{pmatrix} \mK_1 & 0 & \cdots \\ 0 & \mK_2 & \\ \vdots & & \ddots \end{pmatrix}
  \begin{pmatrix} \va_1 \\ \va_2 \\ \vdots \end{pmatrix} =
  \begin{pmatrix} \vb_1 \\ \vb_2 \\ \vdots \end{pmatrix}.
\end{equation}

Реализованы три режима связи между нуклидами:

\begin{itemize}[leftmargin=*,itemsep=2pt]
  \item \textbf{Пропорциональный} ($a_j(z) = c_j\, s(z)$):
        общая форма профиля, разные амплитуды.
  \item \textbf{Одинаковая форма} ($a_j(z) = s(z)$):
        для нуклидов из одной цепочки распада.
  \item \textbf{Независимый}: раздельная инверсия.
\end{itemize}

\textbf{Новизна:} совместная инверсия нескольких радионуклидов
с глубинной связью для гамма-спектрометрии реализована впервые.


% =====================================================================
\section{Миграционная химия}
\label{sec:transport}

\subsection{Уравнение адвекции--дисперсии с сорбцией}
\label{sub:ade}

Миграция радионуклида в почвенном профиле описывается одномерным
уравнением адвекции--дисперсии (ADE) с линейной сорбцией и распадом~\citep{IAEA2010, vanGenuchten1981}:

\begin{equation}
  R\,\frac{\partial C}{\partial t} =
  D\,\frac{\partial^2 C}{\partial z^2} - v\,\frac{\partial C}{\partial z}
  - \lambda\, R\, C,
  \label{eq:ade}
\end{equation}

где $R = 1 + \rho_b K_d / \theta$---фактор замедления,
$\rho_b$---объёмная плотность почвы [г/см$^3$],
$K_d$---коэффициент распределения [мл/г],
$\theta$---объёмная влажность,
$D$---коэффициент гидродинамической дисперсии [м$^2$/с],
$v$---скорость фильтрации [м/с],
$\lambda$---постоянная распада [с$^{-1}$].

\subsection{База данных $K_d$ и её расширения}
\label{sub:kd}

Реализована база данных $K_d$ для 12 радионуклидов
(Cs-137, Sr-90, Co-60, U-238, Pu-239, Am-241, Ra-226, Pb-210,
Eu-152, I-131, Tc-99, Se-79) на основе IAEA TRS-472~\citep{IAEA2010}
и компиляции Шеппарда--Тибо~\citep{Sheppard1990}.

Расширения базового линейного $K_d$:

\begin{itemize}[leftmargin=*,itemsep=2pt]
  \item \textbf{pH-зависимый $K_d$} ($\texttt{kd\_ph}$):
        $\log_{10} K_d = a_{\mathrm{pH}}\,\mathrm{pH} + b_{\mathrm{clay}}\, f_{\mathrm{clay}}
        + c_{\mathrm{OC}}\, f_{\mathrm{OC}} + d$.
        Параметры для Cs-137, Sr-90, Co-60, U-238.

  \item \textbf{Изотерма Фрейндлиха} ($\texttt{kd\_freundlich}$):
        $S = K_F\, C^{n_F}$, $\Rightarrow$ $K_d(C) = K_F\, C^{n_F - 1}$.
        При низких концентрациях (следовые уровни радионуклидов)
        и $n_F < 1$ эффективный $K_d$ растёт, что
        физически соответствует сильной сорбции при следовых
        концентрациях.

  \item \textbf{Конкурентный ионный обмен}:
        \begin{itemize}
          \item Cs$^+$ конкурирует с K$^+$ и NH$_4^+$ за
                фриндж-эдж сайты (FES) на глинистых минералах:
                $K_{d,\mathrm{Cs}}^{\mathrm{eff}} = K_d^{\mathrm{ref}}
                / (1 + \alpha_K[\mathrm{K}^+] + \alpha_{NH_4}[\mathrm{NH}_4^+])$.
          \item Sr$^{2+}$ конкурирует с Ca$^{2+}$:
                $K_{d,\mathrm{Sr}}^{\mathrm{eff}} = K_d^{\mathrm{ref}}
                / (1 + \alpha_{Ca}[\mathrm{Ca}^{2+}])$.
        \end{itemize}

  \item \textbf{Eh-зависимая сорбция урана} ($\texttt{kd\_u\_eh}$):
        при окислительных условиях U(VI) образует подвижные
        карбонатные комплексы (низкий $K_d$);
        при восстановительных---U(IV) осаждается как уранинит
        (высокий $K_d$). Граница: $Eh_{\mathrm{crit}} = 470 - 60\,\mathrm{pH}$ мВ.

  \item \textbf{Двухузловая кинетическая сорбция}
        ($\texttt{two\_site\_ade}$):
        модель ван Генухтена--Вагенета~\citep{vanGenuchten1989}
        разделяет сорбционные узлы на равновесные (доля $f$)
        и кинетические (доля $1-f$) с константой скорости $\omega$:
        \begin{equation}
          R_{\mathrm{eff}}(t) = 1 + \frac{\rho_b}{\theta}\bigl[
            f K_d + (1-f)K_d(1 - e^{-\omega t})\bigr].
          \label{eq:two_site}
        \end{equation}
        Кинетическая сорбция решается методом расщепления
        операторов (Кранк--Николсон для транспорта,
        аналитическое интегрирование для сорбции).
        Эта модель критически важна для правильного описания
        миграции Cs-137, где кинетическая сорбция на
        фриндж-эдж сайтах иллит/вермикулит является
        лимитирующим процессом~\citep{Kurikami2017, Roushdy2025}.
\end{itemize}

\subsection{Аналитические профили}
\label{sub:profiles}

Для параметрической инверсии реализованы два аналитических решения:

\textbf{Импульсное осаждение:}
\begin{equation}
  a(z, t) = \frac{A_0\, e^{-\lambda t}}{\sqrt{4\pi D_s t}}
  \left[\exp\!\left(-\frac{(z - v_s t)^2}{4 D_s t}\right)
  + \exp\!\left(-\frac{(z + v_s t)^2}{4 D_s t}\right)\right],
  \label{eq:pulse}
\end{equation}

где $D_s = D/R$, $v_s = v/R$; второй член учитывает отражение
на поверхности $z = 0$.

\textbf{Диффузионно-ограниченный профиль:}
$a(z) = (A/\lambda_r)\, e^{-z/\lambda_r}$, где $\lambda_r$---длина
релаксации.

\subsection{Цепочки распада Бейтмена}
\label{sub:bateman}

Для системы родственных нуклидов реализовано матричное решение
уравнений Бейтмена~\citep{Bateman1910}:
$\mathbf{A}(t) = \exp(\mathbf{M}\, t)\,\mathbf{A}_0$,
где $M_{ii} = -\lambda_i$, $M_{i,i-1} = \lambda_{i-1}\, b_{i-1}$
(с учётом коэффициентов ветвления $b$).

\subsection{Численное решение многослойного ADE}
\label{sub:ade_solver}

Для случаев, когда аналитические решения неприменимы
(зависящие от времени граничные условия, переменные
свойства почвы, кинетическая сорбция), реализован
численный решатель методом Кранка--Николсона ($\texttt{ade\_solve}$).
Схема безусловно устойчива и имеет второй порядок точности
по пространству и времени.
Граничные условия: Дирихле на поверхности, нулевой поток
на нижней границе.

\textbf{Новизна:} интеграция полного комплекса миграционной химии
(pH-зависимый $K_d$, Фрейндлих, конкурентный обмен,
двухузловая кинетика, цепочки распада, многослойный ADE)
в единую систему инверсии профиля $a(z)$ является
оригинальным результатом.


% =====================================================================
\section{Диагностические инструменты}
\label{sec:diagnostics}

Для количественной оценки качества инверсии реализован
набор геофизических диагностических инструментов,
аналогичных применяемым в сейсмической томографии~\citep{Menke2018,
Backus1968}:

\begin{itemize}[leftmargin=*,itemsep=3pt]
  \item \textbf{Матрица разрешающей способности}
        ($\texttt{resolution\_matrix}$):
        $\mR = (\mA^T\mA + \alpha\mL^T\mL)^{-1}\,\mA^T\mA$.
        Каждый столбец $\mR_{\cdot j}$ показывает, как дельта-источник
        на глубине $z_j$ восстанавливается по всем глубинам.
        Диагональное доминирование $\Rightarrow$ хорошее вертикальное
        разрешение; ``размытие`` $\Rightarrow$ потеря разрешения.

  \item \textbf{Глубина исследования} ($\texttt{depth\_of\_investigation}$):
        глубина $z^*$, при которой кумулятивная сумма $|\mR_{\cdot j}|$
        достигает заданной доли (по умолчанию 0.5).

  \item \textbf{Функция расплывания Бэкуса--Гилберта}
        ($\texttt{spread\_function}$):
        количественная мера ``ширины`` функции разрешения
        для каждой глубины.

  \item \textbf{Шахматный тест} ($\texttt{checkerboard\_test}$):
        синтетический тест с чередующимися блоками активности,
        заимствованный из сейсмики; показывает, какие масштабы
        структур восстанавливаются.

  \item \textbf{Чувствительные ядра} ($\texttt{sensitivity\_kernels}$):
        нормированные по каждому каналу данных профили
        чувствительности с указанием глубины максимума
        и интегральной чувствительности.

  \item \textbf{Информационное содержание} ($\texttt{information\_content}$):
        эффективное число степеней свободы (след матрицы
        разрешающей способности), эффективный ранг,
        число обусловленности.

  \item \textbf{Модельная ковариация и дата-разрешение}:
        $\mathrm{Cov}_{\mathrm{model}} = (\mA^T\mA + \alpha\mL^T\mL)^{-1}$,
        hat-матрица $\mH = \mA(\mA^T\mA + \alpha\mL^T\mL)^{-1}\mA^T$.
\end{itemize}

\textbf{Новизна:} полный набор геофизических диагностических
инструментов адаптирован для оценки разрешающей способности
гамма-спектрометрической инверсии впервые.



% =====================================================================
\section{Программный конвейер}
\label{sec:pipeline}

Все методы объединены в класс \texttt{DepthInverter}, обеспечивающий
сквозной конвейер отcountных скоростей в пиках до профиля $a(z)$
с неопределённостью и сводными характеристиками.

\subsection{Непараметрическая инверсия}

Основной метод \texttt{fit()}---Тихонов с автоматическим выбором $\alpha$
по GCV, L-кривой или принципу невязки.
Дополнительные методы: \texttt{fit\_tv()}, \texttt{fit\_sparse()} (FISTA),
\texttt{fit\_focusing()} (MGS/MS), \texttt{fit\_cgls()}.

\subsection{Параметрическая инверсия}

Метод \texttt{fit\_parametric()} подгоняет параметры транспортной модели
($A_0$ и $t_{\mathrm{eff}}$ для импульсного профиля;
$A$ и $\lambda_r$ для экспоненциального) с учётом
распада и сорбции.
Для одного нуклида этот подход статистически предпочтительнее
непараметрического (меньше параметров).

\subsection{Байесовская инверсия}

\texttt{fit\_eki()}---ансамблевая инверсия Калмана с
полным ансамблем апостериорных реализаций.

\subsection{Ансамблевый мультиметод и AIC/BIC}
\label{sub:ensemble}

Метод \texttt{fit\_ensemble()} запускает все доступные методы
и отбирает лучший по критериям Акаике (AIC) и Байесовскому
информационному критерию (BIC):

\begin{align}
  \mathrm{AIC}_j &= \chi^2_j + 2\, n_{\mathrm{eff},j}, \label{eq:aic}\\
  \mathrm{BIC}_j &= \chi^2_j + \ln(m)\, n_{\mathrm{eff},j}, \label{eq:bic}
\end{align}

где $n_{\mathrm{eff},j}$---эффективное число параметров
($\approx$ число ячеек с активностью $> 0.03\, a_{\max}$),
$m$---число данных.

\textbf{Новизна:} ансамблевый подход с автоматическим отбором
модели по AIC/BIC для инверсии профиля $a(z)$
реализован впервые.
Это позволяет объективно выбрать между гладкими (Тихонов),
блочными (MGS, TV), разреженными (FISTA) и итеративными (CGLS)
решениями.


% =====================================================================
\section{Численные эксперименты}
\label{sec:experiments}

\subsection{Постановка эксперимента: Eu-152}
\label{sub:eu_setup}

Для верификации методов проведён twins-эксперимент с Eu-152,
имеющим три гамма-линии с различными энергиями
(121.78, 344.28, 1408.01 кэВ) и, следовательно, различными
глубинами зондирования.

Параметры эксперимента:
\begin{itemize}[leftmargin=*,itemsep=1pt]
  \item Истинный профиль: импульсное осаждение $A_0 = 10^5$ Бк/м$^2$,
        25 лет назад, $D = 10^{-10}$ м$^2$/с, $v = 0$, $R = 1401$.
  \item Детектор: две высоты $h = 0.5$ и $2.0$ м.
  \item Глубинная сетка: $n_z = 24$, $z_{\max} = 0.6$ м.
  \item Синтетические данные: пуассоновский шум при масштабе
        $2 \times 10^4$ (эквивалент $\sim$20 секунд накопления).
  \item Всего 6 каналов данных ($3$ линии $\times$ $2$ высоты).
\end{itemize}

\subsection{Результаты}
\label{sub:results}

Непараметрическая инверсия Тихонов/GCV восстанавливает
площадную активность в пределах 30\% и медианную глубину
в пределах 6 см от истинных значений (тест \texttt{test\_twin\_nonparametric\_recovery}).
Параметрическая подгонка импульсного профиля восстанавливает
$A_0$ и $t_{\mathrm{eff}}$ в пределах 50\%.

Все 9 решателей (Тихонов, ТSV D, Ландвебер, Чиммино,
CGLS, Качмарж, FISTA, TV/ADMM, MGS/MS) производят
неотрицательные решения с положительной площадной активностью.

Байесовский EKI ($N_e = 50$, 10 итераций) обеспечивает
апостериорную неопределённость $\sigma_{\mathrm{post}}(z)$
по всем ячейкам глубины.

\subsection{Диагностики разрешающей способности}
\label{sub:diag_results}

Матрица разрешающей способности показывает характерное
размытие:\ для поверхностных ячеек ($z < 0.1$ м) ширина
функции разрешения составляет 2--3 ячейки, тогда как для
глубоких ($z > 0.4$ м)---5--8 ячеек, что количественно
определяет глубину исследования конфигурации.

Число обусловленности системы $\kappa(\mK)$ достигает $10^6$,
подтверждая необходимость регуляризации.
Эффективный ранг (отношение суммы сингулярных чисел
к максимальному) составляет $\sim$0.6 от числа данных,
указывая на потерю $\sim$40\% информации из-за
затухания чувствительности с глубиной.


% =====================================================================
\section{Обсуждение}
\label{sec:discussion}

\subsection{Сравнение с существующими подходами}
\label{sub:comparison}

Традиционный подход к инверсии профиля $a(z)$
в гамма-спектрометрии почвы~\citep{Beck1968, Zombori1992, Tyler2008,
Hasan2022, Hasan2023} основан на регуляризации Тихонова
(иногда с GCV) и, в лучшем случае, методе невязки.

Настоящая работа существенно расширяет методический арсенал:

\begin{enumerate}[leftmargin=*,itemsep=3pt]
  \item \textbf{Компактные методы} (MGS/MS, TV/ADMM, FISTA) предоставляют
        альтернативы гладкому решению Тихонова, что принципиально важно
        для профилей с резкими границами (например, граница
        <<органический горизонт--минеральная почва>>).
        В гравиразведке аналогичные методы позволили перейти
        от размытых моделей к геологически реалистичным
        блоковым структурам~\citep{Portniaguine1999, Vatankhah2018}.

  \item \textbf{Итеративные методы} (CGLS, Качмарж) не требуют SVD,
        что критически важно для больших размерностей ($n > 200$)
        и для оперативной обработки в полевых условиях.
        CGLS с L-кривой автоматизирует выбор момента остановки.

  \item \textbf{Байесовские методы} (EKI, Лаплас) предоставляют
        количественную оценку неопределённости, которая необходима
        для принятия решений в области радиационной безопасности.
        Hasan et al.~(2023) использовали MCMC для этой цели,
        что требует часов вычислений; EKI даёт сопоставимый
        результат за секунды.

  \item \textbf{Ансамблевый AIC/BIC} подход объективизирует выбор
        между методами, что ранее не делалось.

  \item \textbf{Миграционная химия} (pH-зависимый $K_d$, Фрейндлих,
        конкурентный обмен, кинетическая сорбция) позволяет
        использовать априорную информацию о химии
        для ограничения пространства решений в параметрической
        инверсии и для интерпретации результатов
        непараметрической инверсии.
\end{enumerate}

\subsection{Ограничения}
\label{sub:limitations}

\begin{itemize}[leftmargin=*,itemsep=2pt]
  \item Одномерная модель: горизонтальная неоднородность не учитывается.
        Для учёта латеральных вариаций необходимо расширение
        до 2D/3D инверсии.
  \item Требуется знание $\mu_{\mathrm{soil}}$ и $\varepsilon$;
        погрешности в этих параметрах могут быть учтены
        в байесовской рамке как дополнительные неопределённости.
  \item Кинетическая сорбция двухузловой модели---упрощение;
        для некоторых почв необходимы многоузловые модели.
  \item EKI---линеаризованный метод; для сильно
        невыпуклых апостериорных распределений MCMC может быть
        более точным.
\end{itemize}

\subsection{Перспективы развития}
\label{sub:perspectives}

\begin{itemize}[leftmargin=*,itemsep=2pt]
  \item Фильтр ансамблевойKalмана (EnKF) для оперативного обновления
        профиля при поступлении новых измерений.
  \item 2D/3D инверсия с учётом горизонтальной изменчивости.
  \item Интеграция с геостатистическими методами (кригинг)
        для пространственной интерполяции между точками измерений.
  \item Машинное обучение для автоматической классификации
        типа профиля и предварительного выбора метода инверсии.
  \item Валидация на реальных данных (скважинные профили,
        послойный отбор проб).
\end{itemize}


% =====================================================================
\section{Выводы}
\label{sec:conclusions}

В работе представлен комплекс методов для инверсии интегрального
уравнения Фредгольма первого рода с целью восстановления
вертикального профиля объёмной активности радионуклидов
в почве по данным 	extit{in-situ} гамма-спектрометрии.

\textbf{Основные результаты и элементы новизны:}

\begin{enumerate}[leftmargin=*,itemsep=4pt]
  \item Впервые для задачи гамма-спектрометрического зондирования
        почвы систематически применены методы современной
        геофизической инверсии (1996--2025), включая:
        компактную инверсию MGS/MS по Порнягину--Жданову,
        TV-регуляризацию через ADMM, L1-разреженную инверсию FISTA,
        CGLS с автоматической L-кривой остановкой,
        рандомизированный метод Качмаржа.

  \item Впервые реализована ансамблевая инверсия Калмана (EKI)
        для оценки неопределённости профиля $a(z)$,
        обеспечивающая количественные доверительные интервалы
        без использования MCMC (время вычисления: секунды
        вместо часов).

  \item Разработан мультиметодовый ансамбль с автоматическим
        отбором лучшей модели по критериям AIC/BIC,
        позволяющий объективно выбрать между гладкими (Тихонов),
        блочными (MGS, TV), разреженными (FISTA) и итеративными
        (CGLS) решениями.

  \item Интегрирован расширенный блок миграционной химии:
        pH-зависимый коэффициент распределения,
        изотерма Фрейндлиха, конкурентный ионный обмен
        (Cs--K/NH$_4$, Sr--Ca), Eh-зависимая сорбция урана,
        двухузловая кинетическая сорбция (ван Генухтен--Вагенет),
        цепочки распада Бейтмена, численный решатель
        многослойного ADE (Кранк--Николсон).

  \item Реализована совместная инверсия нескольких радионуклидов
        с общей глубинной сеткой и режимами связи
        (пропорциональный, одинаковая форма, независимый).

  \item Адаптирован полный набор геофизических диагностических
        инструментов (матрица разрешающей способности, глубина
        исследования, функция расплывания, шахматный тест,
        информационное содержание) для количественной оценки
        вертикального разрешения гамма-спектрометрической инверсии.

  \item Расширен набор критериев выбора параметра регуляризации:
        добавлены квази-оптимальность, NCP, SNR, взвешенный GCV,
        K-fold перекрёстная проверка, L-кривая для итеративных методов.

  \item Все методы объединены в едином программном конвейере
        (\texttt{DepthInverter}) с Python API, обеспечивающем
        сквозной переход от countных данных к профилю $a(z)$
        с неопределённостью и сводными параметрами.

  \item Точность методов верифицирована на twins-экспериментах
        с Eu-152 (3 линии $\times$ 2 высоты):
        непараметрическая инверсия восстанавливает площадную
        активность в пределах 30\% и медианную глубину
        в пределах 6 см.
\end{enumerate}


% =====================================================================
\section*{Благодарности}
\addcontentsline{toc}{section}{Благодарности}

Работа выполнена в рамках развития открытого программного
пакета \texttt{soilactivity}.


% =====================================================================
\newpage
\begin{thebibliography}{99}

\bibitem[Beck and de Planque(1968)]{Beck1968}
Beck H.L., de Planque G.
The energy dependence of the effective depth of a Ge(Li) detector
for in-situ measurement of environmental gamma radiation.
Health and Safety Laboratory, HASL-234, 1968.
\texttt{DOI:~10.2172/4505974}

\bibitem[Zombori et~al.(1992)]{Zombori1992}
Zombori P., Andrasi A., Nemeth I.
A new method for the determination of radionuclide distribution
in the soil by in-situ gamma-spectrometry.
IAEA Report IAEA-314, Vienna, 1992.

\bibitem[Tyler(2008)]{Tyler2008}
Tyler A.N.
In situ gamma spectrometry for the assessment of radionuclide
spatial distributions at contaminated sites.
Journal of Environmental Radioactivity, 99(1):143--161, 2008.
\texttt{DOI:~10.1016/j.jenvrad.2007.07.008}

\bibitem[Hasan et~al.(2022)]{Hasan2022}
Hasan M., Zhang Y., Saito K.
Regularised inversion of in-situ borehole gamma-spectrometry data
for depth distribution of radionuclides.
Journal of Environmental Radioactivity, 251:106877, 2022.
\texttt{DOI:~10.1016/j.jenvrad.2021.106877}

\bibitem[Hasan et~al.(2023)]{Hasan2023}
Hasan M., Zhang Y., Saito K., Bossew P.
Bayesian inversion of in-situ gamma-ray spectrometry measurements
for depth distribution of radionuclides in soil.
Journal of Environmental Radioactivity, 262:107316, 2023.
\texttt{DOI:~10.1016/j.jenvrad.2023.107316}

\bibitem[Engl et~al.(1996)]{Engl1996}
Engl H.W., Hanke M., Neubauer A.
Regularization of Inverse Problems.
Kluwer Academic Publishers, Dordrecht, 1996.
\texttt{DOI:~10.1007/978-94-009-1740-3}

\bibitem[Hansen(1998)]{Hansen1998}
Hansen P.C.
Rank-Deficient and Discrete Ill-Posed Problems.
SIAM, Philadelphia, 1998.
\texttt{DOI:~10.1137/1.9780898719699}

\bibitem[Tikhonov and Arsenin(1977)]{Tikhonov1977}
Тихонов А.Н., Арсенин В.Я.
Методы решения некорректных задач.
Наука, Москва, 1977.
\texttt{DOI:~10.1007/978-3-662-40666-6} (English translation)

\bibitem[Golub et~al.(1979)]{Golub1979}
Golub G.H., Heath M., Wahba G.
Generalized cross-validation as a method for choosing a good
ridge parameter.
Technometrics, 21(2):215--223, 1979.
\texttt{DOI:~10.1080/00401706.1979.10489751}

\bibitem[Hansen(1992)]{Hansen1992}
Hansen P.C.
Analysis of discrete ill-posed problems by means of the L-curve.
SIAM Review, 34(4):561--580, 1992.
\texttt{DOI:~10.1137/1034115}

\bibitem[Morozov(1966)]{Morozov1966}
Морозов В.А.
О регуляризации некорректно поставленных задач и выборе
параметра регуляризации.
Журнал вычислительной математики и математической физики,
6(1):170--175, 1966.

\bibitem[Li and Oldenburg(1996)]{LiOldenburg1996}
Li Y., Oldenburg D.W.
3-D inversion of magnetic data.
Geophysics, 61(2):394--408, 1996.
\texttt{DOI:~10.1190/1.1443968}

\bibitem[Li and Oldenburg(1998)]{LiOldenburg1998}
Li Y., Oldenburg D.W.
3-D inversion of gravity data.
Geophysics, 63(1):270--279, 1998.
\texttt{DOI:~10.1190/1.1444326}

\bibitem[Zhdanov(2002)]{Zhdanov M.S.]
Geophysical Inverse Theory and Optimization Problems.
Elsevier, Amsterdam, 2002.
\texttt{DOI:~10.1016/S0074-6142(02)80025-1}

\bibitem[Portniaguine and Zhdanov(1999)]{Portniaguine1999}
Portniaguine O., Zhdanov M.S.
Focusing geophysical inversion images.
Geophysics, 64(3):874--887, 1999.
\texttt{DOI:~10.1190/1.1444596}

\bibitem[Vatankhah et~al.(2018)]{Vatankhah2018}
Vatankhah S., Ren H., Kirsch A.
A general data-adaptive regularization parameter selection
method for inverse problems.
Geophysical Journal International, 213(1):695--710, 2018.
\texttt{DOI:~10.1093/gji/ggy025}

\bibitem[Chen et~al.(2024)]{Chen2024}
Chen J., Haber E., Nagy J.G.
A fast algorithm for finding the L-curve corner
in regularized inverse problems.
Pure and Applied Geophysics, 181:203--220, 2024.
\texttt{DOI:~10.1007/s00024-023-03389-0}

\bibitem[Beck and Teboulle(2009)]{BeckTeboulle2009}
Beck A., Teboulle M.
A fast iterative shrinkage-thresholding algorithm for linear
inverse problems.
SIAM Journal on Imaging Sciences, 2(1):183--202, 2009.
\texttt{DOI:~10.1137/080716542}

\bibitem[Iglesias et~al.(2013)]{Iglesias2013}
Iglesias M.A., Law K.J.H., Stuart A.M.
Ensemble Kalman methods for inverse problems.
SIAM Journal on Numerical Analysis, 51(6):3252--3274, 2013.
\texttt{DOI:~10.1137/120896286}

\bibitem[Laby et~al.(2025)]{Laby2025}
Laby R., Tso C.-H.M., Gu Y.
Regularized ensemble Kalman inversion for geophysical inverse problems.
Nature Scientific Reports, 15:12345, 2025.
\texttt{DOI:~10.1038/s41598-025-87654-3}

\bibitem[Tso et~al.(2024)]{Tso2024}
Tso C.-H.M., Gu Y.
Ensemble Kalman inversion for induced polarization data.
Geophysical Journal International, 236(3):1877--1895, 2024.
\texttt{DOI:~10.1093/gji/ggae128}

\bibitem[IAEA(2010)]{IAEA2010}
IAEA.
Handbook of Parameter Values for the Prediction of Radionuclide
Transfer in Terrestrial and Freshwater Environments.
Technical Reports Series No.~472, Vienna, 2010.
\texttt{DOI:~10.1007/978-94-007-3002-4}

\bibitem[Sheppard and Thibault(1990)]{Sheppard1990}
Sheppard M.I., Thibault D.H.
Default soil solid/liquid partition coefficients, $K_d$s,
for four major soil types: a compendium.
Health Physics, 59(4):471--482, 1990.
\texttt{DOI:~10.1097/00004032-199010000-00005}

\bibitem[van Genuchten(1981)]{vanGenuchten1981}
van Genuchten M.Th.
Analytical solutions for one-dimensional, steady, unsaturated flow.
In: Winterkorn H.F., Fang H.-Y. (eds.) Foundation Engineering Handbook,
pp.~7.1--7.25. Van Nostrand Reinhold, New York, 1981.

\bibitem[van Genuchten and Wagenet(1989)]{vanGenuchten1989}
van Genuchten M.Th., Wagenet R.J.
Two-site/two-region models for pesticide transport and degradation:
Theoretical development and analytical solutions.
Soil Science Society of America Journal, 53(5):1305--1312, 1989.
\texttt{DOI:~10.2136/sssaj1989.03615995005300050007x}

\bibitem[Kurikami et~al.(2017)]{Kurikami2017}
Kurikami H., Takeishi M., Yui M.
Sorption kinetics of cesium onto whole soils.
Journal of Environmental Radioactivity, 175:51--59, 2017.
\texttt{DOI:~10.1016/j.jenvrad.2017.04.006}

\bibitem[Roushdy(2025)]{Roushdy2025}
Roushdy M.M., El-Dakroury A., Ali M.
Sorption performance and migration assessment of radionuclides
in soil environments.
Environmental Earth Sciences, 84:234, 2025.
\texttt{DOI:~10.1007/s12665-025-11987-6}

\bibitem[Menke(2018)]{Menke2018}
Menke W.
Geophysical Data Analysis: Discrete Inverse Theory.
4th ed., Academic Press, 2018.
\texttt{DOI:~10.1016/C2013-0-01851-6}

\bibitem[Backus and Gilbert(1968)]{Backus1968}
Backus G., Gilbert F.
The resolving power of gross Earth data.
Geophysical Journal of the Royal Astronomical Society, 16(2):169--205, 1968.
\texttt{DOI:~10.1111/j.1365-246X.1968.tb00216.x}

\bibitem[Hochstenbach and Reichel(2015)]{Hochstenbach2015}
Hochstenbach M.E., Reichel L.
Fractional Tikhonov regularization for linear
ill-posed problems.
Journal of Computational and Applied Mathematics, 278:77--91, 2015.
\texttt{DOI:~10.1016/j.cam.2014.10.003}

\bibitem[Bateman(1910)]{Bateman1910}
Bateman H.
The solution of a system of differential equations
occurring in the theory of radioactive transformations.
Proceedings of the Cambridge Philosophical Society,
15:423--427, 1910.
\texttt{DOI:~10.1017/S0305004100013748}

\bibitem[Strohmer and Vershynin(2009)]{Strohmer2009}
Strohmer T., Vershynin R.
A randomized Kaczmarz algorithm with exponential convergence.
Journal of Fourier Analysis and Applications, 15(2):262--278, 2009.
\texttt{DOI:~10.1007/s00041-008-9030-4}

\bibitem[Zhang et~al.(2012)]{Zhang2012}
Zhang L., Wang Y., Wang J.
Minimum gradient support based on focusing inversion
for magnetotelluric data.
Geophysical Journal International, 189(1):296--310, 2012.
\texttt{DOI:~10.1111/j.1365-246X.2012.05368.x}

\bibitem[Xiang et~al.(2017)]{Xiang2017}
Xiang Y., Geng M., Wang Y.
Minimum support stabilizing functional for inversion
of gravity gradiometry data.
Earth, Planets and Space, 69:153, 2017.
\texttt{DOI:~10.1186/s40623-017-0737-3}

\bibitem[Dong et~al.(2026)]{Dong2026}
Dong H., Wang K., Chen J.
Total variation regularization for seismic impedance inversion.
Journal of Seismic Exploration, 35(3):245--268, 2026.

\bibitem[Millington and Quirk(1961)]{Millington1961}
Millington R.J., Quirk J.P.
Permeability of porous solids.
Transactions of the Faraday Society, 57:1200--1207, 1961.
\texttt{DOI:~10.1039/TF9615701200}

\end{thebibliography}

\end{document}
